\SourceLeaf{96}%

\noindent Es handelt sich also um eines aus der Reihe jener Theoreme, die \textsc{Fourier} seit 1807 fand und 1822 in der »Théorie analytique de la chaleur« ausführlich veröffentlichte. Die von \textsc{Gauss} zum Beweis des Theorems angegebene Summenformel liefert für $\operatorname{lim.}\,\omega=0$ unmittelbar die Integralformel von \textsc{Fourier}. Dabei muss die Gleichung $\mathop{\mathrm{lim.}}\limits_{x=\pm\infty}\varphi(x)=0$ gelten. Diese Bedingung ist für die von \textsc{Gauss} in der Wahrscheinlichkeitsrechnung verwendete Funktion $\varphi(x)$ tatsächlich erfüllt (vgl. »Theoria combinationis observationum etc.« Art. 4, Gesammelte Werke IV S. 5).

\Needspace{17\baselineskip}
Was die Summenformel von \textsc{Gauss} selbst betrifft, so lässt sie sich zunächst durch Aufteilen des Integrals auf der rechten Seite in die folgende Form bringen.
\[
\sum^{+\infty}_{k=-\infty}\varphi\!\left(t+\frac{2\pi k}{\omega}\right) = \frac{\omega}{2\pi}\sum^{+\infty}_{l=-\infty}\left(e^{-ilt\omega}\mathord{.}\sum^{+\infty}_{k=-\infty}\int^{\frac{2(k+1)\pi}{\omega}}_{\frac{2k\pi}{\omega}}\varphi(\theta)\mathord{.} e^{il\theta\omega}\,d\theta\right).
\]
\Needspace{15\baselineskip}
\noindent Führt man hier eine neue Integrationsvariable $x$ durch $\theta\omega = x+2k\pi$ ein und vertauscht auf der rechten Seite die Reihenfolge der Summationen, so erhält man die folgende neue Form der Summenformel.
\[
\sum^{\infty}_{k=-\infty}\varphi\!\left(t+\frac{2\pi k}{\omega}\right) = \sum^{+\infty}_{k=-\infty}\sum^{+\infty}_{l=-\infty}\left(e^{-ilt\omega}\mathord{.}\frac{1}{2\pi}\int^{2\pi}_{0}\varphi\!\left(\frac{x+2k\pi}{\omega}\right)e^{ilx}\,dx\right).
\]
\Needspace{12\baselineskip}
\noindent Die Richtigkeit dieser Formel folgt unmittelbar aus der Reihe von \textsc{Fourier}.
\[
\varphi\!\left(t+\frac{2\pi k}{\omega}\right) = \sum^{+\infty}_{l=-\infty}\left(e^{-ilt\omega}\mathord{.}\frac{1}{2\pi}\int^{2\pi}_{0}\varphi\!\left(\frac{x+2k\pi}{\omega}\right)e^{ilx}\,dx\right).
\]

\begin{flushright}
\textsc{Fricke}.
\end{flushright}
